\title{DevSecOps 파이프라인을 활용한 단계적 Hybrid PQC-TLS 전환 자동화}
\titlerunning{단계적 Hybrid PQC-TLS 전환 자동화}

\author{
First Author\inst{1}\orcidID{0000-1111-2222-3333} \and
Second Author\inst{1}\orcidID{1111-2222-3333-4444} \and
Third Author\inst{1}\orcidID{2222--3333-4444-5555} \and
Fourth Author\inst{1}\orcidID{3333-4444-5555-6666} \and
Hwa-Jeong Seo\inst{1}\orcidID{0000-0003-0069-9061}
}
%
\authorrunning{F. Author et al.}
%
\institute{Hansung University, Seoul, Korea\\
\email{hwajeong84@gmail.com}}
%
\maketitle
\begin{abstract}
양자 컴퓨터 실용화 이전에 암호화 통신을 수집하여 이후 해독하는 ``지금 수집, 나중에 해독(Harvest Now, Decrypt Later, HNDL)'' 공격은 TLS 통신 구간에 장기적 위협을 제기한다.
NIST PQC 표준화 완료 이후 단계적 TLS 전환이 요구되고 있으나, 기존 방식은 설정 오류의 뒤늦은 발견, 암호 자산 추적 부재, 자동 대응 미비라는 한계를 가진다.
본 논문은 Policy as Code 기반 11-Step DevSecOps 파이프라인을 제안한다.
제안 시스템은 Classical ECC(Stage~1), Hybrid PQC(Stage~2), Advanced Hybrid PQC(Stage~3)의 3단계 전환을 자동 검증·배포하며, OQS-Nginx 자체 빌드, 보안 스캔, CBOM Diff, 매트릭스 CI 회귀 감지를 통합한다.
평가 결과, PQC 도구 체인 함정 6종에 대해 표준 SCA/SAST/SBOM은 0/6을 탐지한 반면 자체 빌드는 5/6 예방, 매트릭스 CI는 6/6을 감지하였다.
LLM 기반 코드 마이그레이션 80 trial 평가에서 정적 AST 검증만으로는 런타임 실패를 차단할 수 없으며, 의미적 실행 검증이 필수적임을 확인하였다.

\keywords{양자 내성 암호 \and DevSecOps \and TLS 전환 \and Hybrid PQC-TLS \and CI/CD 파이프라인}

\end{abstract}